\chapter{A-tól Z-ig: egy minimális webalkalmazás}

\noindent\textbf{Tanulási út: Gyors út: alap.} Készítsd el az első teljes, végponttól végpontig tartó vertikális szeletet.\par\medskip

\rustbridgeblockhu{A függvények, \code{Result}, \code{?} hibaterjesztés, lokális változók és hétköznapi vezérlési szerkezetek ugyanazt a szerepet töltik be, mint Rustban.}{A callable attribútumok, route-ok, typed request mezők, HTML-építés, context és capability köti ezt a számítást a HTTP-hez.}{Ne találj ki külön Velran-specifikus hibaterjesztési vagy lokálisváltozó-stílust; a Rust-szerű mag körüli webes határra koncentrálj.}


\invariantblock{Biztonsági invariáns: explicit route policy}{Route soha nem lesz implicit public. Minden HTTP belépési pont deklarálja a hozzáféréshez szükséges authorityt; a hiányzó policy hiba, nem alapértelmezés.}

\diagnosticblockhu{Compiler diagnostic}{SEC-A01-001}{Egy route nem kapott explicit hozzáférési policyt. A Velran nem találgatja, hogy a végpont public vagy authenticated legyen.}{Deklaráld a szándékolt policyt a route-on, például \code{public} vagy \code{auth user}; az alkalmazás contractja alapján válassz, ne csak a diagnostic elnémítására.}

Mielőtt a nyelvi részleteket külön tanulnánk, építsünk egy teljes, szándékosan kicsi alkalmazást. Csak minimális development konfigurációt használ, és megmutatja az utat a forrástól az adatbázison és HTTP-n át a telepítésig; nem kész webshop.

A példa egy mini terméklista. A kezdőoldal kilistázza az adatbázisban lévő termékeket, és ugyanott egy formmal új terméket is fel lehet venni. Ezzel egyetlen példában megjelenik a model, typed SQL, GET page, POST action, form validation, CSRF, transaction, migration, konfiguráció és a reverse proxy mögötti futtatás.

A repositoryban a kész példa itt található:

\begin{lstlisting}
examples/a-z-products/
  main.vrn
  migrations/0001_create_products.sql
  server.toml.example
  README.md
\end{lstlisting}

\section{A request életciklusa: egyetlen mentális modell a legtöbb webes feladathoz}

A könyv alkalmazáskódjának nagy része ugyanazzal a hétlépcsős modellel ellenőrizhető:

\begin{center}
\begin{tabularx}{\textwidth}{@{}lY@{}}
\toprule
Lépés & Kérdés \\
\midrule
Request & Melyik route és HTTP metódus fogadta a kérést? \\
Parse & Mely deklarált query, form, JSON, upload vagy path értékből lesz típusos input? \\
Validate & Milyen bound, refinement vagy több mezőt érintő szabály kötelező? \\
Authenticate & Ki a hívó, és szükséges-e session vagy gépi identity? \\
Authorize & A hívó végrehajthatja-e ezt a műveletet ezen a konkrét objektumon vagy tenanton? \\
Query/tranzakció & Milyen olvasás vagy állapotmódosítás engedett az előző kapuk után? \\
Response & Milyen típusos HTML, JSON, redirect, státusz vagy hiba hagyja el az alkalmazást? \\
\bottomrule
\end{tabularx}
\end{center}

A lépések gondolkodási sorrendet jelentenek, nem azt, hogy minden request hét külön runtime függvényt futtat. Egy publikus GET-hez nem feltétlenül kell autentikáció; egy read-only oldal nem feltétlenül nyit tranzakciót. A fegyelem lényege, hogy minden kérdést explicit felteszünk, és csak akkor hagyunk ki egy lépést, ha a contract szerint valóban nincs rá szükség.

Állapotmódosító kérésnél őrizd meg a security sorrendet: protected munka előtt parse és validate, identityfüggő döntés előtt authenticate, protected effect előtt authorize, majd a módosítás a megfelelő tranzakciós vagy critical-operation contract alatt történjen. A response csak ezeken a határokon átjutott adatból épüljön fel.

Hibakereséskor ugyanezen a láncon haladj balról jobbra. Azt keresd, hol tér el először a tényleges viselkedés a deklarált contracttól, ahelyett hogy a ``handlert'' egyetlen átláthatatlan lépésként kezelnéd.

\section{A könyvtár létrehozása}

Fejlesztéshez legyen egy külön alkalmazáskönyvtárunk:

\begin{lstlisting}
mkdir -p my-products/migrations
cd my-products
\end{lstlisting}

Ebben a példában minden Velran kód egyetlen \code{main.vrn} fájlban marad. Nagyobb alkalmazásnál később \code{mod} fájlokra bontjuk.

\section{Az adatbázis-séma}

A \code{migrations/0001_create_products.sql} tartalma SQLite esetén:

\begin{lstlisting}[language=SQL]
CREATE TABLE products (
    id INTEGER PRIMARY KEY AUTOINCREMENT,
    name TEXT NOT NULL,
    price INTEGER NOT NULL CHECK (price >= 0)
);
\end{lstlisting}

Ez a migration hozza létre azt a táblát, amelyet az alkalmazás typed SQL lekérdezései használnak. Productionban a migráció külön deploy lépés; a server nem módosítja automatikusan a sémát induláskor.

\section{A model és a lekérdezések}

A \code{main.vrn} elején a modellközpontú fejezetben megismert módon először a typed üzleti adatalakot írjuk le:

\begin{lstlisting}[language=Velran]
model Product {
    id: i64
    name: String
    price: i64
}
\end{lstlisting}

A lista lekérdezése:

\begin{lstlisting}[language=Velran]
#[query]
fn listProducts(db: Db) -> Result<Vec<Product>, DbError> sql {
    SELECT id, name, price
    FROM products
    ORDER BY id DESC
    LIMIT 50
}
\end{lstlisting}

A \code{Vec<Product>} azt mondja ki, hogy a query nulla vagy több \code{Product} sort adhat vissza. A \code{SELECT} oszlopai pontosan a model mezőinek nevét és sorrendjét követik.

Az INSERT transaction capabilityt kap:

\begin{lstlisting}[language=Velran]
#[query]
fn createProduct(
    tx: Transaction,
    name: String,
    price: i64
) -> Result<(), DbError> sql {
    INSERT INTO products(name, price)
    VALUES (:name, :price)
}
\end{lstlisting}

A \code{:name} és \code{:price} bind paraméter. Nem szövegkonkatenációval építünk SQL-t felhasználói inputból.

\section{A GET oldal}

A kezdőoldal lekéri a termékeket, megjelenít egy formot és kilistázza az eredményt:

\begin{lstlisting}[language=Velran]
#[page]
fn home(ctx: PageContext, db: Db) -> Result<Html, PageError> {
    let products = listProducts(db)?;

    return Ok(html {
        <main>
            <h1>Mini termeklista</h1>

            <form method="post" @action(create)>
                <input type="hidden" name="_csrf" value="{{ csrfToken }}">
                <label>
                    Nev
                    <input name="name" maxlength="100">
                </label>
                <label>
                    Ar
                    <input name="price">
                </label>
                <button type="submit">Hozzaadas</button>
            </form>

            <ul>
                @for product in products {
                    <li>
                        #{{ product.id }} - {{ product.name }} - {{ product.price }}
                    </li>
                }
            </ul>
        </main>
    });
}
\end{lstlisting}

A dinamikus szövegek HTML-escape-elve kerülnek a válaszba. A form action nem kézzel írt URL, hanem az \code{@action(create)} helperből származik. A rejtett \code{_csrf} mező a POST kérés CSRF-védelméhez kell.

\section{A POST action}

A form sikeres beküldése transactionben ír az adatbázisba, majd redirectel vissza a kezdőoldalra:

\begin{lstlisting}[language=Velran]
#[action]
fn create(
    ctx: ActionContext,
    db: Db,
    name: String,
    price: i64
) -> Result<Redirect, PageError> {
    transaction db {
        createProduct(tx, name, price)?;
    }
    return Ok(redirect(home()));
}
\end{lstlisting}

A redirect célja a typed GET route helper, \code{home()}, nem string URL. Ez a Post/Redirect/Get minta: ha a felhasználó frissíti a visszakapott oldalt, a böngésző nem küldi be újra automatikusan az INSERT-et.

\section{A route-ok és a validáció}

A fájl végén két route köti össze a HTTP felületet a handlerekkel:

\begin{lstlisting}[language=Velran]
route home GET "/" public => home;

route create POST "/products"
    form name<String> price<i64>
    validate name length 1 100 price range 0 100000000
    public => create;
\end{lstlisting}

A POST route már a handler előtt typed inputtá alakítja a két mezőt. Üres vagy 100 karakternél hosszabb név, negatív ár, túl nagy ár vagy nem egész szám esetén a kérés nem jut el érvényes üzleti műveletként az INSERT-ig.

A publikus írás ebben a kicsi tanpéldában szándékos egyszerűsítés. Valódi public submission endpointnál explicit döntsd el, kell-e authentication, route rate limit, abuse monitoring vagy erősebb business permission. Ne másold át a \code{public} policyt csak azért, mert a példa ezt használja.

\section{A teljes alkalmazás egy helyen}

Az előző részek együtt alkotják a teljes alkalmazást: \code{model} + két \callable{query} + \callable{page} + \callable{action} + két \code{route}. Nem ismételjük meg még egyszer ugyanazt a hosszú listinget; a repositoryban a pontos, egyfájlos változat közvetlenül megnyitható:

\begin{lstlisting}
examples/a-z-products/main.vrn
\end{lstlisting}

Ez szándékos könyvszerkesztési döntés is: a következő fejezetekben az egyes nyelvi elemeket már részletesen tárgyaljuk, itt viszont az alkalmazás teljes útvonalán akarunk végigmenni.

\section{SQLite kapcsolat developmenthez}

A database URL-t ne írjuk a forrásba. Készítsünk külön fájlt, abszolút SQLite pathszal:

\begin{lstlisting}
printf '%s\n' 'sqlite:///tmp/velran-a-z-products.db' > dev-db-url
chmod 600 dev-db-url
\end{lstlisting}

A fájl itt nem valódi titkot tartalmaz, de ugyanazt a \code{*_file} mintát használjuk, mint productionben a jelszavas database URL esetén.

\section{Forrásellenőrzés és migration}

Ha a release binárisokat már telepítettük például \code{/usr/local/bin} alá:

\begin{lstlisting}
/usr/local/bin/velran-cli check main.vrn

/usr/local/bin/velran-cli migrate verify \
  --dir migrations \
  --db-url-file "$PWD/dev-db-url"

/usr/local/bin/velran-cli migrate apply \
  --dir migrations \
  --db-url-file "$PWD/dev-db-url"
\end{lstlisting}

A \code{check} a nyelvi és statikus szerződést ellenőrzi. A migration CLI külön folyamat; ettől még az alkalmazásszerver nem indult el.

\section{Minimális development konfiguráció}

A \code{server.toml} legyen például:

\begin{lstlisting}
[server]
app = "/ABS/PATH/my-products/main.vrn"
listen = "127.0.0.1:8080"
insecure_dev_cookies = true

[database]
url_file = "/ABS/PATH/my-products/dev-db-url"
\end{lstlisting}

A két \code{/ABS/PATH} helyére valódi abszolút útvonal kell. Az \code{insecure_dev_cookies} csak loopback developmenthez elfogadható; productionben ne kapcsold be.

Először ellenőrizzük a konfigurációt:

\begin{lstlisting}
/usr/local/bin/velran-server --config "$PWD/server.toml" --check-config
\end{lstlisting}

Majd indulhat a szerver:

\begin{lstlisting}
/usr/local/bin/velran-server --config "$PWD/server.toml"
\end{lstlisting}

\section{Az első HTTP kérés}

Másik terminálból:

\begin{lstlisting}
curl -i http://127.0.0.1:8080/
\end{lstlisting}

Böngészőben nyisd meg ugyanazt az URL-t. A form beküldése után a folyamat:

\begin{enumerate}
  \item a böngésző POST-ot küld a \code{/products} route-ra;
  \item a runtime ellenőrzi a request schema-t, a typed mezőket, a validációt és a CSRF tokent;
  \item a \code{create} action transactiont nyit;
  \item a \code{createProduct} bind paraméterekkel INSERT-et futtat;
  \item siker esetén commit történik;
  \item a handler redirectet ad a \code{/} URL-re;
  \item az új GET már az adatbázisból visszaolvasott sort rendereli.
\end{enumerate}

Ez az egyszerű lánc a későbbi wiki és CMS alkalmazásban is ugyanaz marad; csak több model, query, route, jogosultság és komponens kerül köré.

\section{Mit történik hibás inputnál?}

Próbálj negatív árat vagy nem numerikus értéket megadni. A fontos elv az, hogy a böngésző HTML attribútumai csak felhasználói kényelmet adnak; a valódi szerződést a szerveroldali route schema és validation adja. Az SQL query csak már typed értéket kap.

Érdemes még a következőket kipróbálni:

\begin{itemize}
  \item küldj hiányzó \code{name} mezőt;
  \item küldj ismeretlen extra form mezőt;
  \item módosítsd a migrationt alkalmazás után, majd futtasd újra a \code{migrate verify} parancsot;
  \item állítsd le az adatbázis elérését, és figyeld meg a readiness illetve a szerverlog viselkedését.
\end{itemize}

\section{Developmentből production felé}

A példa alkalmazáskódja ugyanaz marad release környezetben is; nem kell újraírni csak azért, mert a szervert már telepített binárissal futtatjuk. A \ref{ch:telepites}. fejezetben bemutatott build és telepítés után a szerver tipikusan így indul:

\begin{lstlisting}
/usr/local/bin/velran-server \\
  --config /usr/local/etc/velran/server.toml
\end{lstlisting}

Az alkalmazás entrypointja lehet például \code{/srv/velran/current/main.vrn}. Apache vagy Nginx esetén a publikus TLS kapcsolatot a proxy terminálja, a Velran pedig loopbacken hallgat. A teljes Apache/Nginx mintát nem ismételjük meg itt: a telepítési fejezet mutatja be a trust boundaryt, a könyv végi függelék pedig copy/paste-közeli production konfigurációt ad.

\section{Mit tanultunk ebből az egy példából?}

Egy mindössze néhány tucat soros alkalmazásban már láttuk a Velran teljes alap útvonalát:

\begin{lstlisting}
HTTP GET
  -> route
  -> page
  -> typed query
  -> database
  -> escaped HTML

HTTP POST
  -> route schema + validation + CSRF
  -> action
  -> transaction
  -> typed SQL bind
  -> commit
  -> redirect
  -> GET
\end{lstlisting}

A következő fejezetek ezt a képet bontják szét. Amikor egy-egy nyelvi vagy runtime részlet technikainak tűnik, érdemes visszagondolni erre a két kérésre: \emph{hol lép be az adat a rendszerbe, és milyen typed/capability határokon halad át, mielőtt tartós állapot vagy HTTP válasz lesz belőle?}

\section{Ellenőrzött companion példa}
A fejezet teljes, compilerrel ellenőrzött companion példája az \path{examples/a-z-products/main.vrn}. A könyv kódblokkjai vezetett részletek; teljes program másolásakor a repository fájlt tekintsd elsődlegesnek, mert a \code{./verify.sh} pontosan ezt a forrást ellenőrzi.

\section{Gyakorlat: kövess végig egy kérést}
\paragraph{Gyakorlási mód: vezetett.} Használd az előszóban meghatározott vezetett módszert.

Válassz egy GET és egy POST kérést a példaalkalmazásból, és mondd végig az összes határt: route-illesztés, típusos input-dekódolás, autorizáció, query vagy tranzakció végrehajtása, válaszépítés és redirect. Ha egy határt nem tudsz megnevezni, térj vissza ahhoz a részhez, amíg az adatfolyam egyértelmű nem lesz.

\section{Tipikus hiba: összekeverjük az olvasási és módosítási folyamatot}
Egy GET page és egy POST action érintheti ugyanazt a domain objektumot, de más biztonsági és életciklus-szemantikájuk van. Az olvasás maradjon mellékhatásmentes; a módosítás menjen validáción, autorizáción, tranzakciós állapotváltozáson és ahol szükséges Post/Redirect/Get (PRG) mintával.

\section{Folyamatos mintaprojekt: Secure Notes}
Valósítsd meg a legkisebb hasznos vertikális szeletet: a jelenlegi felhasználó jegyzeteinek listázását és egy új jegyzet létrehozását. Még ne adj hozzá szerkesztést, megosztást, keresést, feltöltést vagy cache-t. A cél a teljes request--storage--response útvonal bizonyítása a lehető legkevesebb mozgó résszel.
